\scanpage{67}
% archive sheet number: 68
\margin{\unsure{à mettre au net|à mettre au point|à mettre après} Weyl \unsure{sur les propriétés} dans les \unsure{Hilbert|hilbertiens}}

% margin note written sideways
\margin{à \unsure{développer} en rédigeant le cas de la dimension finie}

% margin note written sideways
\margin{\deleted{\uline{\unsure{Conséquences|Conséquence}, \unsure{changement de notation}.} Pour la suite, les \unsure{notations} suivantes seront utiles. Si $u$ est un op. compact dans un espace de Hilbert complexe, on désigne par $(\lambda_i(u))$ la suite des valeurs propres de $u$ \unsure{répétées|ordonnées|comptées} suivant leurs multiplicités, \deleted{\unsure{de}} rangées de façon que \unsure{la suite en module|leur suite en module} soit décroissante. Si \unsure{cette liste} ne \unsure{contient} qu'un nombre fini de valeurs propres non nulles, on complète par suite \deleted{\unsure{d'une infinité de valeurs nulles}}, en posant \unsure{$\lambda_i(u)=0$ pour les autres $i$}. Si $u$ est un op. compact \illegible{2}, \unsure{$(\rho_i(u)=(\lambda_i((u^*u)^{1/2}))$|$\rho_i(u)=(\lambda_i(uu^*u)^{1/2})$}, suite \inserted{décroissante} des valeurs propres de $(u^*u)^{1/2}$ (qui est la même que celle des valeurs \unsure{propres} de $(uu^*)^{1/2}$). \deleted{\unsure{La définition doit s'étendre}} \unsure{Si $u$ est unitaire} \illegible{2}. On obtient \unsure{ainsi tout naturellement}, en \unsure{liaison} avec les notations. Si \unsure{donc} $u$ est un op. compact positif d'un \unsure{espace}, $\lambda_i(u)=\rho_i(u)=\unsure{\sigma_i(u)}$.}}

\deleted{\unsure{Supposons}}

\uline{Proposition 2} \deleted{\unsure{Si $u$ est un}} Soit $u$ un op. compact dans l'espace hilbertien complexe \deleted{\unsure{H}}$E$, $(\lambda_i)$ la suite de ses valeurs propres. \deleted{\unsure{comprises}} Alors \deleted{les valeurs} $\overset{n}{\Lambda}u$ est un op. compact dans $\overset{n}{\Lambda}^{(2)}E$, ayant \inserted{exactement} le \unsure{système} de valeurs propres \deleted{\unsure{les valeurs|propres}} $\lambda_{i_1}\ldots\lambda_{i_n}$, avec $i_1<i_2<\ldots<i_n$.

Comme l'application $\deleted{\unsure{E|F}}u\longrightarrow\overset{n}{\Lambda}u$ \deleted{est une \unsure{appl.}} de $L(E)$ dans $L(\overset{n}{\Lambda}^{\unsure{(2)|(1)}}E)$ est continue, \deleted{si $u$ est compact, il est La question de \unsure{rang f|rang fini}} et transforme les opérateurs de rang fini en opérateurs de rang fini, il s'ensuit que si $u$ est compact, donc limite d'opérateurs de rang fini, \deleted{il en est de même} \inserted{alors} $\overset{n}{\Lambda}u$ est aussi limite d'opérateurs de rang fini, donc aussi un op. compact.

D'autre part, si $u$ est de rang fini, \deleted{il \unsure{peut} se mettre sous forme d'un} \deleted{\inserted{qui à une base hilb., si $u$ \unsure{commutait|convenable}} \unsure{certain} \unsure{espace} \unsure{tel} dont un nb fini seulement, \inserted{\unsure{Alors} soit $N$, un \unsure{tel}, considérons} soit $E_1$ l'espace} soit $E_{1}$ l'intersection de son noyau $\unsure{u^{-1}(0)|u^{-1}0}$ avec l'orthogonal de \unsure{l'image} $u(E)$, et soit $E_0$ l'orthogonal de $E_1$.

Alors $u$ \uline{s'annule} sur $E_1$, et $u(E)\subset E_0$, donc $u$ \uline{\unsure{conserve|invarie}} $E_0$. $\overset{n}{\Lambda}E$ est \deleted{la somme \unsure{di}} \inserted{algébriquement} isomorphe à la somme $\sum_{0\leq p\leq n}\overset{n-p}{\Lambda}E_0\otimes\overset{p}{\Lambda}E_1\deleted{\unsure{\}}}$, et quand on munit d'une \deleted{\unsure{$E$|$F$|$\Lambda E$}} part $\overset{n}{\Lambda}E$ de sa structure préhilbertienne, et d'autre part $\overset{n-p}{\Lambda}E_0$ et $\overset{p}{\Lambda}E_1$ de leur structure préhilbertienne naturelle, donc sur $\overset{n-p}{\Lambda}E_0\otimes\overset{p}{\Lambda}E_1$ une structure préhilbertienne naturelle, on vérifie que l'isomorphisme $\overset{n}{\Lambda}E=\sum_{\unsure{0\leq p\leq n}}\overset{n-p}{\Lambda}E_0\otimes\overset{p}{\Lambda}E_1$
